\input{build/foundations-preamble.tex}
\usepackage[authoryear]{natbib}
\usepackage{bussproofs}
\renewenvironment{prooftree}
  {\begin{center}\bottomAlignProof}
  {\DisplayProof\end{center}}
\usepackage[tableaux]{prooftrees}
\newenvironment{oltableau}{\center\tableau{}}{\endtableau\endcenter}
\renewcommand{\refname}{સંદર્ભગ્રંથો}
\renewcommand{\partname}{ભાગ}
\newenvironment{intro}{}{}
\newenvironment{derivation}{%
  ~\begin{trivlist}\item\begin{tabular}[b]{@{}rll@{}}}
  {\end{tabular}\end{trivlist}}
\newcommand{\Id}[1]{\mathord{\mathrm{Id}_{#1}}}
\newcommand{\equivrep}[2]{[#1]_{#2}}
\newcommand{\equivclass}[2]{#1/_{\!{#2}}}
\newcommand{\Intequiv}{\mathrel{\sim_{\Int}}}
\newcommand{\Ratequiv}{\mathrel{\sim_{\Rat}}}
\newcommand{\Realequiv}{\mathrel{\sim_{\Real}}}
\newcommand{\funrestrictionto}[2]{#1\mathord{\restriction}_{#2}}
\newcommand{\funimage}[2]{#1[#2]}
\newcommand{\emptyseq}{\Lambda}
\newcommand{\liff}{\mathbin{\leftrightarrow}}
\newcommand{\dom}[1]{\mathrm{dom}(#1)}
\newcommand{\ran}[1]{\mathrm{ran}(#1)}
\newcommand{\comp}[2]{#2\circ #1}
\newcommand{\pto}{\mathrel{\ooalign{\hfil$\mapstochar\mkern5mu$\hfil\cr$\to$}}}
\newcommand{\fdefined}{\downarrow}
\newcommand{\fundefined}{\uparrow}
\newcommand{\True}{\mathbb{T}}
\newcommand{\False}{\mathbb{F}}
\newcommand{\lfalse}{\bot}
\newcommand{\ltrue}{\top}
\newcommand{\Obj}[1]{\mathsf{#1}}
\newcommand{\Lang}[1]{\mathcal{#1}}
\NewDocumentCommand{\Frm}{o}{%
  \IfNoValueTF{#1}{\mathrm{Frm}}{\mathrm{Frm}(\Lang{#1})}}
\newcommand{\PVar}{\mathrm{At}_0}
\newcommand{\Struct}[1]{\mathfrak{#1}}
\newcommand{\pAssign}[1]{\mathfrak{#1}}
\newcommand{\pValue}[1]{\overline{\mathfrak{#1}}}
\NewDocumentCommand{\pSat}{t{/} m m}{%
  \pAssign{#2}\IfBooleanTF{#1}{\nvDash}{\vDash}#3}
\newcommand{\Entails}{\vDash}
\newcommand{\Proves}{\vdash}
\newcommand{\ident}{\equiv}
\newcommand{\subst}[2]{#1/#2}
\newcommand{\SSubst}[2]{#1[#2]}
\newcommand{\Subst}[3]{#1[\subst{#2}{#3}]}
\newcommand{\indfrm}{}
\newcommand{\indcase}[3]{\def\indfrm{#1}$#1\ident #2$: #3}
\newcommand{\Sequent}{\Rightarrow}
\renewcommand{\fCenter}{\ensuremath{\,\Sequent\,}}
\newcommand{\LeftR}[1]{\ensuremath{{#1}\mathrm{L}}}
\newcommand{\RightR}[1]{\ensuremath{{#1}\mathrm{R}}}
\newcommand{\Weakening}{\ensuremath{\mathrm{W}}}
\NewDocumentCommand{\Intro}{m o}{%
  \ensuremath{{#1}\mathrm{Intro}\IfNoValueTF{#2}{}{_{#2}}}}
\NewDocumentCommand{\Elim}{m o}{%
  \ensuremath{{#1}\mathrm{Elim}\IfNoValueTF{#2}{}{_{#2}}}}
\newcommand{\FalseInt}{\ensuremath{\lfalse_I}}
\NewDocumentCommand{\Discharge}{m m}{[#1]^{#2}}
\NewDocumentCommand{\DischargeRule}{m m}{%
  \RightLabel{#1}\LeftLabel{\scriptsize $#2$}}
\NewDocumentCommand{\sFmla}{m m o}{%
  \ensuremath{\IfNoValueTF{#3}{}{#3\,}\hbox to.8em{\ensuremath{#1}\hfil}#2}}
\NewDocumentCommand{\TRule}{m m o}{%
  \ensuremath{{#2}{#1}\IfNoValueTF{#3}{}{\,#3}}}
\newcommand{\TAss}{ધારણા}
\expandafter\def\csname guformula@A\endcsname{\varphi}
\expandafter\def\csname guformula@B\endcsname{\psi}
\newcommand{\formula}[1]{\csname guformula@#1\endcsname}
\expandafter\def\csname gutoken@formula\endcsname{સૂત્ર}
\expandafter\def\csname gutoken@valuation\endcsname{સત્યમૂલ્ય-નિયુક્તિ}
\expandafter\def\csname gutoken@derivation\endcsname{નિષ્પત્તિ}
\expandafter\def\csname gutoken@derivation@P\endcsname{નિષ્પત્તિઓ}
\expandafter\def\csname gutoken@tableau\endcsname{ટેબ્લો}
\expandafter\def\csname gutoken@tableau@P\endcsname{ટેબ્લો}
\renewcommand{\usetoken}[2]{%
  \ifcsname gutoken@#2@#1\endcsname
    \csname gutoken@#2@#1\endcsname
  \else
    \csname gutoken@#2\endcsname
  \fi}
\let\printtoken\usetoken
\newcommand{\sourcecorrection}[2]{%
  \leavevmode\par\smallskip\noindent
  \fbox{\begin{minipage}{0.94\linewidth}\small
    \textbf{સ્રોત-સુધારો #1.} #2
  \end{minipage}}%
  \par\smallskip}
\definecolor{oldiagcolorD}{HTML}{1973ba}
\newlength{\olphotowidth}
\setlength{\olphotowidth}{.4\textwidth}
\RenewDocumentCommand{\olasset}{O{\olphotowidth} m}{\centering\resizebox{#1}{!}{\input{upstream/#2}}}
\input{build/proof-systems-available-labels.tex}
\renewcommand{\oliflabeldef}[3]{\ifcsname guavailable@#1\endcsname#2\else#3\fi}
\let\guoriginalref\ref
\expandafter\def\csname guexternalref@sfr:siz::chap\endcsname{\guoriginalref{sfr:siz:int:sec}}
\expandafter\def\csname guexternalref@sfr:arith::chap\endcsname{\guoriginalref{sfr:arith:int:sec}}
\expandafter\def\csname guexternalref@his:set:mythology:sec\endcsname{%
  \href{https://github.com/OpenLogicProject/OpenLogic/blob/9620cc73f9c8e0ad003c514a5d3748f29611c4c0/content/history/set-theory/mythology.tex}{મૂળ ગ્રંથનો સંબંધિત ઐતિહાસિક વિભાગ}}
\expandafter\def\csname guexternalref@his:set:limits:sec\endcsname{%
  \href{https://github.com/OpenLogicProject/OpenLogic/blob/9620cc73f9c8e0ad003c514a5d3748f29611c4c0/content/history/set-theory/limits.tex}{મૂળ ગ્રંથનો લક્ષોનો વિભાગ}}
\expandafter\def\csname guexternalref@sth:::part\endcsname{%
  \href{https://github.com/OpenLogicProject/OpenLogic/tree/9620cc73f9c8e0ad003c514a5d3748f29611c4c0/content/set-theory}{મૂળ ગ્રંથનો ગણસિદ્ધાંત ભાગ}}
\expandafter\def\csname guexternalref@sth:ord-arithmetic::chap\endcsname{%
  \href{https://github.com/OpenLogicProject/OpenLogic/tree/9620cc73f9c8e0ad003c514a5d3748f29611c4c0/content/set-theory/ord-arithmetic}{મૂળ ગ્રંથનું ક્રમસંખ્યાઓના અંકગણિતનું પ્રકરણ}}
\renewcommand{\ref}[1]{%
  \ifcsname guexternalref@#1\endcsname
    \csname guexternalref@#1\endcsname
  \else
    \guoriginalref{#1}%
  \fi}
\hypersetup{pdftitle={ગણો, વિધાનાત્મક તર્કશાસ્ત્ર અને નિષ્પત્તિ તંત્રો — ઓપન લોજિક ગુજરાતી}}
\begin{document}
\renewcommand{\ref}[1]{%
  \ifcsname guexternalref@#1\endcsname
    \csname guexternalref@#1\endcsname
  \else
    \guoriginalref{#1}%
  \fi}
\begin{titlepage}
{\LARGE\bfseries ગણો, સંબંધો, વિધેયો, ગણોનું કદ, અંકગણિતીકરણ,
અનંત ગણો, વિધાનાત્મક તર્કશાસ્ત્ર અને નિષ્પત્તિ તંત્રો\par}
\vspace{1em}
{\Large ઓપન લોજિક — ગુજરાતી\par}
\vspace{2em}
ઓપન લોજિક પ્રોજેક્ટના અંગ્રેજી મૂળનો ગુજરાતી અનુવાદ.\par
\vspace{1em}
આ આવૃત્તિમાં અગાઉના પ્રકરણો સાથે નિષ્પત્તિ તંત્રોના સામાન્ય
પરિચય અને ચાર મુખ્ય પદ્ધતિઓના સર્વેક્ષણના પાંચ ખંડ ઉમેર્યા છે.
તેમાં સિક્વન્ટ કલન, સ્વાભાવિક નિગમન, ટેબ્લો અને સ્વયંસિદ્ધિમૂલક
નિષ્પત્તિ સાથે યથાર્થતા, પૂર્ણતા અને સુસંગતતા સામેલ છે.\par
\vfill
આ યંત્ર દ્વારા કરેલો અનુવાદ છે. સંપૂર્ણ ગ્રંથનું કામ ચાલુ છે.
આ આવૃત્તિમાં મૂળનાં ૬૫ એકમોનો સમાવેશ છે; કુલ ૭૨૨ એકમો છે.
કેટલીક વિશિષ્ટ પરિભાષા કામચલાઉ છે. ચકાસણીની વિગતો અને
મૂળ પાઠ વિશેની નોંધો સંપાદકીય માહિતીમાં આપેલી છે.\par
\vspace{1em}
\href{https://github.com/OpenLogicProject/OpenLogic}{Open Logic Project}
\quad \href{https://creativecommons.org/licenses/by/4.0/}{CC BY 4.0}\par
\href{https://github.com/KokunoYumeto/OpenLogic-translations}{અનુવાદોનું કેન્દ્ર}\par
\end{titlepage}
\tableofcontents
\clearpage
\input{build/proof-systems-body.tex}
\clearpage
\section*{સંપાદકીય નોંધો}
\addcontentsline{toc}{section}{સંપાદકીય નોંધો}
આ નોંધો મૂળ પાઠથી અલગ છે. મૂળની ગણિતીય નિશાનીઓ જાળવી છે.
OLFUN-001થી OLFUN-005, OLSIZ-001થી OLSIZ-012, OLARI-001થી
OLARI-004, OLINF-001થી OLINF-002, OLPL-001થી OLPL-002 અને
OLPRF-001થી OLPRF-004 સુધીની ઓળખવાળી નોંધો સ્થિર અંગ્રેજી મૂળમાં
મળેલી ખામીઓ અને ગુજરાતી પાઠમાં કરેલા મર્યાદિત સુધારા બાજુબાજુ દર્શાવે
છે. દરેક સુધારાની ચોક્કસ જગ્યા, કારણ અને સૂત્રમાં થયેલો ફેરફાર
નિર્ણય-અભિલેખમાં નોંધાયેલ છે.

નિષ્પત્તિ તંત્રોના આ સર્વેક્ષણના ખંડો પ્રથમ-ક્રમ અને વિધાનાત્મક
તર્કશાસ્ત્ર બંનેમાં વપરાતા મૂળમાંથી લેવાયા છે. આ વાંચનરૂપમાં સ્થિર
મૂળની વિધાનાત્મક ગોઠવણી પસંદ કરી છે. તેથી સૂત્રો અને નિષ્પત્તિના
દાખલા વિધાનાત્મક તર્કશાસ્ત્ર માટે વાંચવાના છે; સામાન્ય પરિભાષા બંને
સંદર્ભોને લાગુ પડે છે.

સિક્વન્ટની સામાન્ય રજૂઆતમાં જમણી શ્રેણીનો અંતિમ ક્રમસૂચક OLPRF-001
વડે $n$ કર્યો છે. ટેબ્લોના સર્વેક્ષણમાં અસત્ય-સંયોજનના નિયમનું નામ
OLPRF-002 વડે નિયામક નિયમ-સારણી સાથે સરખું કર્યું છે. અસુસંગતતાની
શરતમાં દર્શાવેલી દરેક ધારણા મૂળ ગણની સભ્ય હોવી જરૂરી છે; આ વ્યાપ
OLPRF-003 વડે નિયામક વ્યાખ્યા પ્રમાણે પુનઃસ્થાપિત કર્યો છે. દાખલાના
સત્ય-સંયોજનના બંને નિયમ-નામ OLPRF-004 વડે સુધાર્યાં છે.

નિષ્પત્તિ, સ્વયંસિદ્ધિ, અનુમાન, સુસંગતતા, પૂર્ણતા અને સાબિતી જેવા
પદો ગુજરાતી વિશ્વકોશ અને પાઠ્યસામગ્રીના તપાસેલા પ્રયોગો સાથે
સરખાવ્યાં છે. સિક્વન્ટ કલન, ટેબ્લો અને સાબિતી-સૈદ્ધાંતિક અર્થવિચાર
જેવા વિશિષ્ટ પદો અર્થ સાથે વ્યાખ્યાયિત કર્યા છે અને નિષ્ણાત સુધારા
માટે ખુલ્લા છે.

વિધાનાત્મક તર્કશાસ્ત્રના અગાઉના વાંચનરૂપમાં સ્થિર મૂળની સામાન્ય
ગોઠવણી અનુસરીને અસત્ય અને સત્યના અચળો તથા પાંચેય તાર્કિક સંયોજકોને
મૂળભૂત ગણ્યા છે. વૈકલ્પિક રીતે વ્યાખ્યાયિત સંયોજકોનો અનુવાદ
ગોઠવણી-ટૅગો સાથે સ્રોત ફાઇલમાં જાળવેલો છે.

ગણનીય માટેનું ગુજરાતી પ્રમાણ ગણનીયને ગણનીય અનંત ગણ માટે વાપરે છે,
જ્યારે Open Logicની આ આવૃત્તિની વ્યાખ્યા સાન્ત અને ખાલી ગણોને પણ
સમાવે છે. બંને પરિગણના-રૂપો પોતાની સ્થાનિક વ્યાખ્યા સાથે જ વાંચવાનાં છે.

પરિભાષા, મૂળ સ્રોતો, પ્રત્યેક અનુવાદિત ખંડની સ્રોત-સંગતિ અને
ચકાસણીની વધુ વિગતો આવૃત્તિની સાથેની ફાઇલોમાં છે.
\begin{thebibliography}{9}
\bibitem[Benacerraf(1965)]{Benacerraf1965}
Benacerraf, Paul. 1965. \emph{What numbers could not be}.
The Philosophical Review 74(1), 47--73.
\bibitem[Cantor(1892)]{Cantor1892}
Cantor, Georg. 1892. Über eine elementare Frage der Mannigfaltigkeitslehre.
\bibitem[Frege(1884)]{Frege1884}
Frege, Gottlob. 1884. \emph{Die Grundlagen der Arithmetik}.
\bibitem[Potter(2004)]{Potter2004}
Potter, Michael. 2004. \emph{Set Theory and Its Philosophy}.
\bibitem[Conway(2006)]{Conway2006}
Conway, John. 2006. \emph{The Power of Mathematics}.
\bibitem[Katz and Katz(2012)]{KatzKatz2012}
Katz, Karin Usadi and Mikhail G. Katz. 2012. Stevin Numbers and Reality.
\bibitem[O'Connor and Robertson(2005)]{OConnorRobertson:RN}
O'Connor, John J. and Edmund F. Robertson. 2005. The real numbers: Stevin to Hilbert.
\bibitem[Hilbert(2013)]{EwaldSieg2013}
Hilbert, David. 2013. \emph{David Hilbert's Lectures on the Foundations of Arithmetic and Logic 1917--1933}.
Edited by William Bragg Ewald and Wilfried Sieg. Springer.
\bibitem[Dedekind(1888)]{Dedekind1888}
Dedekind, Richard. 1888. \emph{Was sind und was sollen die Zahlen?} Vieweg.
\end{thebibliography}
\end{document}
